%%%%%%%%%%%%%%%%%%%%%%preface.tex%%%%%%%%%%%%%%%%%%%%%%%%%%%%%%%%%%%%%%%%%
% sample preface
%
% Use this file as a template for your own input.
%
%%%%%%%%%%%%%%%%%%%%%%%% Springer %%%%%%%%%%%%%%%%%%%%%%%%%%

%\preface
\chapter*{Preface}
%% Please write your preface here
In 1982, Richard Feynman pointed out that
a simulation of quantum systems on classical computers
is generally inefficient because the dimension of the state space
increases exponentially with the number of particles~\cite{Feynman82}.
Instead, quantum systems could be simulated efficiently by other quantum systems.
David Deutsch put this idea forward
by formulating a quantum version of 
a Turing machine~\cite{Deutsch85}.
Quantum computation enables us 
to solve certain kinds of problems that 
are thought to be intractable with classical computers
such as the prime factoring problem
%~\cite{Shor} 
and an approximation of the Jones polynomial.
%~\cite{AharonovJones,AharonovJonesHard}.
It has a great possibility to disprove the extended (strong) Church-Turing thesis,
i.e., that any computational process on realistic devices can be simulated 
efficiently on a probabilistic Turing machine.

However, for this statement to make sense,
we need to determine whether or not quantum computation
is a realistic model of computation.
Rolf Landauer criticized it (encouragingly)
by suggesting to put a footnote:
``{\it This proposal, like all proposals for quantum computation, 
relies on speculative technology, 
does not in its current form take into account 
all possible sources of noise, unreliability 
and manufacturing error, and
probably will not work.}"~\cite{LandauerLloyd}.
Actually, quantum coherence,
which is essential for quantum computation
is quite fragile against noise.
%Moreover, quantum information cannot be copied
%due to the no-cloning theorem~\cite{Nocloning}.
If we cannot handle the effect of noise,
quantum computation is of a limiting interest,
like classical analog computers,
as a realistic model of computation. 
To solve this, many researchers have investigated 
the fault-tolerance of quantum computation
with developing quantum error correction techniques.
One of the greatest achievements of this approach 
is topological fault-tolerant quantum computation
using the surface code
%~\cite{Kitaev,Dennis}
proposed by R. Raussendorf {\it et al.}~\cite{RaussendorfAnn,RaussendorfNJP,RaussendorfPRL}.
It says that 
nearest-neighbor two-qubit gates and single-qubit operations
on a two-dimensional array of qubits can perform 
universal quantum computation fault-tolerantly
as long as the error rate per operation is less than $\sim 1\%$.

In this book, 
I present 
a self-consistent review of 
topological fault-tolerant quantum computation
using the surface code.
The book covers 
everything required to 
understand topological 
fault-tolerant quantum computation,
ranging from the definition of the surface code
to topological quantum error correction and 
topological operations on the surface code.
The basic concepts and powerful tools for understanding 
topological fault-tolerant quantum computation, such as 
universal quantum computation, quantum algorithms, stabilizer formalism, 
and measurement-based quantum computation,
are also introduced in the first part (Chapter 1 and Chapter 2) of the book.
In particular, in Chapter 1, I also mention
a quantum algorithm for approximating 
the Jones polynomials,
which is also related to topological quantum computation
with braiding non-Abelian anyons.
In Chapter. 3, the definition of the surface code and 
topological quantum error correction on it is explained. 
In Chapter 4, topological quantum computation on the surface code
is described in the circuit-based model,
where topological diagrams are introduced
to understand the logical operations on the surface code diagrammatically.
In Chapter. 5, 
I explain the same thing in the measurement-based model,
as done in the original proposal~\cite{RaussendorfAnn}.
Hopefully, it would be easy to see how these two viewpoints are related.

Throughout the book, 
I have tried to explain the quantum operations 
using circuit and topological diagrams 
so that the readers can get a graphical understanding of the operations.
The graphical understanding should be helpful to study
the subjects more efficiently.
Topological quantum error correction codes 
are a nice play ground for studying
the interdisciplinary connections between 
quantum information and other fields of physics,
such as condensed matter physics and statistical physics.
Actually, there is a nice correspondence between 
topological quantum error correction codes and topologically ordered
systems in condensed matter physics.
Furthermore, if we consider a decoding problem of 
a quantum error correction code,
a partition function of a random statistical mechanical model
is naturally appeared as a posterior probability for the decoding.
These interdisciplinary topics are also included in Chapter 3.

Almost all topics, except for the basic concepts 
in the first part, are based on the results
achieved after the appearance of the standard textbook
of quantum information science
entitled {\it``Quantum Computation and Quantum Information"} (Cambridge University Press 2000) by M. A. Nielsen and I. L. Chuang.
In this sense, the present comprehensive review 
on these topics would be helpful to learn and update
the recent progress efficiently.
In this book, 
I concentrated on the quantum information aspect 
of topological quantum computation.
Unfortunately, I cannot cover 
the more physical and condensed matter aspects of 
topological quantum computation, such as non-Abelian anyons
and topological quantum field theory.
In this sense, this book is complemented by the book 
{\it``Introduction to topological quantum computation"}
(Cambridge University Press 2012)
written by J. K. Pachos.
Readers who are interested in the more physical aspects
of topological quantum computation are recommended to read it.

Hopefully, this review will encourage 
both theoretical and experimental
researchers to find a more feasible way of quantum computation.
It will also bring me great pleasure if this review provides 
an opportunity 
to reunify and refine
various subdivided fields of modern physics 
in terms of quantum information.







\vspace{\baselineskip}
\begin{flushright}\noindent
Kyoto, Japan
\hfill {\it Keisuke  Fujii}, \\
April, 2015
\\
\end{flushright}


